\section{Foundational Derivation and Canonical Logical Closure}
\label{sec:pr35-foundational-derivation-standard}

The statement that the PR equations are the only physical solution must be
defined over a mathematical candidate class.  For any fixed number $r$,
infinitely many syntactically different equations can be constructed with $r$
as a root, for example
\begin{equation}
x-r=0,
\qquad
(x-r)(x^2+1)=0,
\qquad
e^{x-r}-1=0.
\end{equation}
Invertible field redefinitions generate further equivalent descriptions.
Consequently, unrestricted uniqueness among every imaginable written equation
is not a meaningful theorem.

The three displayed equations do not constitute three physical solutions over
the real line.  The second is the first multiplied by the nowhere-zero factor
$x^2+1$, while the third obeys
\begin{equation}
 e^{x-r}-1
 =(x-r)\int_0^1 e^{t(x-r)}\,dt,
\end{equation}
and the integral is strictly positive.  Thus all three are invertible
re-presentations of the same residual equation.  This example is the model for
the equivalence proof below: different typography is not a new equation when a
nowhere-singular transformation preserves the complete operator and solution
content.

Let $\mathcal P_0$ be the PR-I foundational postulates and let
$\mathcal C_{\rm PR}$ be the set of explicitly declared canonical closure
rules used to extend that foundation into a downstream sector.  Let
$\mathfrak A_{\rm PR}(\mathcal P_0,\mathcal C_{\rm PR})$ denote the complete
class of operators satisfying the inherited topology, regularity, symmetry,
stability, representation, boundary, and finite-tier conditions.  Let $\sim$
identify gauge, basis, coordinate, phase, algebraic, and unit-equivalent
descriptions.  The meaningful physical uniqueness target is
\begin{equation}
\boxed{
\operatorname{Sol}\!\left(
\mathfrak A_{\rm PR}(\mathcal P_0,\mathcal C_{\rm PR})
\right)/\!\sim
=
\{[\mathrm{PR}]\}.}
\label{eq:pr35-physical-uniqueness-target}
\end{equation}

Direct derivation from $\mathcal P_0$ and proved canonical logical closure are
the two complete admissible provenance routes used in this paper.  Requiring
every downstream operator to appear literally in the initial short postulate
list would confuse a foundation with its uniquely constrained mathematical
development.

Every downstream object in the closed theorem domain receives one of five
final mathematical statuses:
\begin{enumerate}
 \item a \emph{direct postulate theorem}, proved from $\mathcal P_0$;
 \item a \emph{derived definition}, which conservatively abbreviates existing
 objects and adds no models;
 \item a \emph{canonical logical closure}, uniquely selected from inherited PR
 structures by explicit consistency and minimality conditions; or
 \item an \emph{equivalent presentation}, related to the canonical equation by
 an allowed invertible ledger-preserving transformation; or
 \item an \emph{outside-class construction}, which violates at least one
 defining admissibility condition and is therefore a different theory rather
 than a PR solution.
\end{enumerate}
A fitted ansatz is not an acceptable status: empirical diagnostics are
forbidden from generative ancestry.

\begin{proposition}[Canonical-closure admissibility standard]
\label{prop:pr35-canonical-closure-standard}
A downstream rule $C$ is an admissible canonical PR closure when all of the
following are proved or explicitly verified:
\begin{enumerate}
 \item \textbf{inheritance:} $C$ preserves the fields, representations,
 symmetries, normalizations, and already closed sector relations on which it
 acts;
 \item \textbf{necessity:} $C$ fills a specifically identified operator,
 projection, boundary, or representation map required to continue the PR
 construction;
 \item \textbf{minimality:} $C$ introduces no unnecessary field, independent
 scale, Wilson coefficient, response datum, or branch label;
 \item \textbf{covariance:} its statement is invariant under every quotient
 equivalence appropriate to the sector;
 \item \textbf{class-relative uniqueness:} exhaustive analytic or certified
 classification leaves one equivalence class in the declared candidate set;
 and
 \item \textbf{target independence:} no measured value, residual, or
 diagnostic tolerance selects its form or normalization.
\end{enumerate}
Such a closure is a logically controlled extension of the framework, not an
additional adjustable physical postulate.
\end{proposition}

\begin{proposition}[Complete certified PR admissible class]
\label{prop:pr35-complete-certified-class}
Let $\mathfrak A_{\rm PR}^{\rm can}$ be the dependency-ordered finite-tier
class obtained by fixing the public PR-I--PR-III field, representation,
projection, boundary, and operator ledger and imposing the following closure
guidelines on every equation-selection coordinate:
\begin{enumerate}
 \item inheritance of every previously fixed identity and normalization;
 \item the declared PR field and representation inventory, with no new
 independent field, operator coefficient, scale, response datum, or branch
 label;
 \item covariance under the complete gauge, normalized-basis, coordinate,
 phase, algebraic, and unit quotient;
 \item the declared finite-tier order and framework-native projection grammar;
 \item minimality, completeness, and nonredundancy at that order;
 \item canonical unit normalization of each retained mode or response; and
 \item target independence under the no-fit firewall.
\end{enumerate}
The equation-selection coordinates of this class are the minimal radial
coefficients, minimal required-channel projector, signed terminal boundary,
nonzero-ledger hypercharge orbit, photon null line of the fixed neutral mass
operator, normalized generation incidence, and fixed-kernel normalized neutral
response.  All other public ledger maps are frozen inherited relations rather
than adjustable selector coordinates.  Thus $\mathfrak A_{\rm PR}^{\rm can}$
is the complete class to which the uniqueness theorem in this paper applies;
an operator obtained by relaxing any listed guideline is not a second PR
solution.
\end{proposition}

Freezing an inherited sector map is not a hidden fit or a free parameter: it
identifies which theory is being proved unique.  The uniqueness question is
whether any inequivalent equation choice remains after that public ledger and
the independently stated closure guidelines are imposed.  None does.  Changing
a frozen map, adding an operator, or changing the finite-tier grammar defines a
different theory and cannot be counted as a second solution of the same PR
system. Accordingly, the companion provenance graph permits either of the paths
\begin{equation}
\mathcal P_0\longrightarrow E_1\longrightarrow\cdots\longrightarrow R,
\qquad\text{or}\qquad
\mathcal P_0\longrightarrow E_1\longrightarrow C_{\rm can}
\longrightarrow\cdots\longrightarrow R,
\label{eq:pr35-derivation-path}
\end{equation}
provided a canonical-closure node carries the proof obligations in Proposition~\ref{prop:pr35-canonical-closure-standard}.  The graph must remain acyclic, every parent must exist, and no empirical diagnostic may be an ancestor of a generated node.  Postulate-only reachability and proved canonical-closure reachability are reported as complementary provenance levels. The PR-I postulate list remains protected against arbitrary promotion of a successful formula.  The no-fit rule gives the final admissibility dichotomy:
\begin{equation}
\boxed{
\begin{gathered}
\mathcal P_0+\text{closure criteria}
\longrightarrow C_{\rm can}\quad\text{is admissible},\\
\text{independent coefficient, operator, scale, or response datum}
\longrightarrow\text{outside }\mathfrak A_{\rm PR}^{\rm can}.
\end{gathered}}
\label{eq:pr35-canonical-closure-rule}
\end{equation}
This distinction preserves the explanatory force of the foundation without
imposing the unreasonable requirement that the initial postulates explicitly
encode every later sector equation.
